%==============================================================================
%  QuarkyLab GPU Cluster — Researcher Guide
%  NetFRAME Home Lab
%  Build:  tectonic QuarkyLab-Researcher-Guide.tex
%  Source-of-truth prose also lives in QuarkyLab-Researcher-Guide.md
%==============================================================================
\documentclass[11pt]{article}

%------------------------------------------------------------------------------
%  CONFIGURATION
%------------------------------------------------------------------------------
\newcommand{\LabName}{QuarkyLab}
\newcommand{\LabHost}{quarkylab}
\newcommand{\AccessHost}{quarkylab.kylemason.org}
\newcommand{\GPUName}{NVIDIA RTX 8000 (48\,GB)}
\newcommand{\HomeQuota}{150\,GB}
\newcommand{\Contact}{our shared Discord channel}
\newcommand{\GuideVersion}{1.0}
\newcommand{\GuideOwner}{NetFRAME Home Lab}
%------------------------------------------------------------------------------

\usepackage[T1]{fontenc}
\usepackage[utf8]{inputenc}
\usepackage{newtxtext,newtxmath}
\usepackage[margin=1in,headheight=15pt]{geometry}
\usepackage{microtype}
\usepackage{xcolor}
\usepackage{enumitem}
\usepackage{booktabs}
\usepackage{array}
\usepackage{tabularx}
\usepackage[most]{tcolorbox}
\usepackage{listings}
\usepackage{titlesec}
\usepackage{fancyhdr}
\usepackage{titletoc}
\usepackage[hidelinks]{hyperref}

%--- palette ------------------------------------------------------------------
\definecolor{quarkA}{HTML}{0B3D66}
\definecolor{quarkB}{HTML}{1E88A8}
\definecolor{quarkAccent}{HTML}{E8A33D}
\definecolor{codebg}{HTML}{F6F8FA}
\definecolor{codeframe}{HTML}{D6DEE7}
\definecolor{codekw}{HTML}{0B3D66}
\definecolor{codecom}{HTML}{6A9A5B}
\definecolor{codestr}{HTML}{A6522C}
\definecolor{noteblue}{HTML}{2C6FB5}
\definecolor{warnred}{HTML}{B23B36}
\definecolor{tipgreen}{HTML}{2E7D46}
\definecolor{rulegray}{HTML}{C9D3DD}

\hypersetup{
  colorlinks=true, linkcolor=quarkA, urlcolor=quarkB, citecolor=quarkB,
  pdftitle={\LabName{} GPU Cluster --- Researcher Guide},
  pdfauthor={\GuideOwner}, bookmarksnumbered=true
}

%--- headers / footers --------------------------------------------------------
\pagestyle{fancy}
\fancyhf{}
\renewcommand{\headrulewidth}{0.4pt}
\renewcommand{\footrulewidth}{0.4pt}
\fancyhead[L]{\small\textsc{\LabName{} Researcher Guide}}
\fancyhead[R]{\small\leftmark}
\fancyfoot[L]{\small \GuideOwner}
\fancyfoot[C]{\small v\GuideVersion}
\fancyfoot[R]{\small\thepage}

%--- section styling ----------------------------------------------------------
\titleformat{\section}{\Large\bfseries\color{quarkA}}{\thesection}{0.6em}{}[{\color{rulegray}\titlerule[0.8pt]}]
\titleformat{\subsection}{\large\bfseries\color{quarkB}}{\thesubsection}{0.6em}{}
\titlespacing*{\section}{0pt}{2.2ex plus 1ex minus .2ex}{1.2ex plus .2ex}

%--- code listings ------------------------------------------------------------
\lstdefinestyle{shell}{
  language=bash,
  basicstyle=\ttfamily\small,
  backgroundcolor=\color{codebg},
  frame=single, framesep=6pt, rulecolor=\color{codeframe},
  xleftmargin=10pt, xrightmargin=6pt,
  keywordstyle=\color{codekw}\bfseries,
  commentstyle=\color{codecom}\itshape,
  stringstyle=\color{codestr},
  showstringspaces=false, breaklines=true, breakatwhitespace=false,
  columns=fullflexible, keepspaces=true, upquote=true,
  escapeinside={(*@}{@*)},
  morekeywords={ssh,ssh-keygen,scp,rsync,cloudflared,winget,brew,apptainer,conda,curl,bash,tmux,screen,nohup,nvidia-smi,python,cat}
}
\lstset{style=shell}
\newcommand{\code}[1]{\texttt{#1}}

%--- admonitions --------------------------------------------------------------
\newtcolorbox{notebox}{enhanced,breakable,colback=noteblue!6,colframe=noteblue,
  coltitle=white,fonttitle=\bfseries,title=Note,boxrule=0.6pt,arc=2pt,left=8pt,right=8pt,top=5pt,bottom=5pt}
\newtcolorbox{tipbox}{enhanced,breakable,colback=tipgreen!7,colframe=tipgreen,
  coltitle=white,fonttitle=\bfseries,title=Tip,boxrule=0.6pt,arc=2pt,left=8pt,right=8pt,top=5pt,bottom=5pt}
\newtcolorbox{warnbox}{enhanced,breakable,colback=warnred!6,colframe=warnred,
  coltitle=white,fonttitle=\bfseries,title=Important,boxrule=0.6pt,arc=2pt,left=8pt,right=8pt,top=5pt,bottom=5pt}

\newcommand{\osbadge}[1]{\vspace{0.4ex}\noindent{\color{quarkA}\rule{2pt}{10pt}}\ \textbf{\color{quarkA}#1}\par\nopagebreak\vspace{0.3ex}}

\setlist[enumerate]{leftmargin=1.6em,itemsep=2pt,topsep=3pt}
\setlist[itemize]{leftmargin=1.4em,itemsep=2pt,topsep=3pt}
\renewcommand{\arraystretch}{1.25}

%==============================================================================
\begin{document}

%--- title block --------------------------------------------------------------
\begin{center}
{\color{quarkA}\rule{\linewidth}{2pt}}\\[0.4cm]
{\huge\bfseries\color{quarkA} \LabName{} GPU Cluster}\\[0.25cm]
{\Large\color{quarkB} Researcher Guide}\\[0.3cm]
{\normalsize Remote access + the commands to run your research}\\[0.35cm]
{\color{quarkA}\rule{\linewidth}{2pt}}\\[0.5cm]
\end{center}

\begin{tcolorbox}[colback=quarkB!6,colframe=quarkB,arc=3pt,boxrule=0.8pt]
\centering\itshape Researchers get a normal interactive shell with direct GPU access --- no SLURM required. This guide assumes you are comfortable at the command line.\\[2pt]
Admin, key requests, and help: \textbf{contact Kyle on \Contact.}
\end{tcolorbox}

\vspace{0.4cm}

%==============================================================================
\section{Access (one-time)}
%==============================================================================
You reach \LabName{} through a \textbf{Cloudflare Tunnel} --- no VPN, no inbound ports, and the server's IP stays hidden. Login is by \textbf{SSH key}. Do this once.

\subsection{Create an SSH key and send the public half}
\begin{lstlisting}
ssh-keygen -t ed25519 -C "you@lab"
cat ~/.ssh/id_ed25519.pub
\end{lstlisting}
Send that \code{.pub} line to Kyle on \Contact. He adds it to your account and tells you your \textbf{username} (e.g.\ \code{cephandrius}). \textbf{Never} send the private key (\code{id\_ed25519}, the one without \code{.pub}).

\subsection{Install cloudflared}
\osbadge{macOS}
\code{brew install cloudflared}

\osbadge{Windows (PowerShell)}
\code{winget install -{}-id Cloudflare.cloudflared}

\osbadge{Linux}
Package or static binary from \url{https://github.com/cloudflare/cloudflared/releases/latest}

\subsection{Route SSH through the tunnel}
Add this to \code{\textasciitilde/.ssh/config} (create the file if it does not exist):
\begin{lstlisting}
Host (*@\texttt{\LabHost}@*)
    HostName (*@\texttt{\AccessHost}@*)
    ProxyCommand cloudflared access ssh --hostname %h
    IdentityFile ~/.ssh/id_ed25519
\end{lstlisting}

\subsection{Connect}
\begin{lstlisting}
ssh <your-username>@(*@\texttt{\LabHost}@*)
\end{lstlisting}
\code{cloudflared} dials out on demand and your key logs you in (accept the host-key prompt with \code{yes} the first time). You land on a normal shell in your home directory.

\begin{notebox}
An email sign-in gate (Cloudflare Access) may be added later; if so, your first connection will also open a one-time browser login. Nothing you set up above changes.
\end{notebox}

%==============================================================================
\section{The machine}
%==============================================================================
\begin{itemize}
  \item 1\,$\times$ \textbf{Quadro RTX 8000, 48\,GB} --- driver 550.163.01, CUDA 12.1 stack.
  \item \code{nvidia-smi} shows GPU status and current usage. You share this \textbf{one physical GPU} with students' batch jobs and the researcher --- glance at it before grabbing a lot of VRAM.
\end{itemize}

%==============================================================================
\section{GPU and the ML environment}
%==============================================================================
The curated stack is an Apptainer image whose \code{ml} conda environment has \textbf{Python 3.11, PyTorch 2.5.1+cu121, TensorFlow 2.21}, plus NumPy, pandas, SciPy, scikit-learn, and the HuggingFace stack. Image: \code{/data/containers/base.sif} (fast local copy at \code{/workspace/containers/base.sif}).
\begin{lstlisting}
# run a script on the GPU
apptainer exec --nv /data/containers/base.sif python train.py

# interactive shell inside the container
apptainer shell --nv /data/containers/base.sif

# quick sanity check
apptainer exec --nv /data/containers/base.sif \
  python -c "import torch; print(torch.cuda.is_available(), torch.cuda.get_device_name(0))"
# -> True Quadro RTX 8000
\end{lstlisting}
\begin{itemize}
  \item \code{python} inside the image is already the \code{ml} env (\code{/opt/conda/envs/ml}).
  \item \code{-{}-nv} is what exposes the GPU --- omit it and CUDA is invisible.
  \item Your \textbf{home} and \textbf{current directory} are auto-mounted; bind extra paths with \code{-B /host/path:/container/path}.
\end{itemize}

%==============================================================================
\section{Your own environment (optional)}
%==============================================================================
No sudo, but your home is yours. For a custom stack, install Miniforge in your home:
\begin{lstlisting}
curl -L -o ~/miniforge.sh https://github.com/conda-forge/miniforge/releases/latest/download/Miniforge3-Linux-x86_64.sh
bash ~/miniforge.sh -b -p ~/miniforge3
~/miniforge3/bin/conda create -n myenv python=3.11 pytorch-cuda=12.1 ...
\end{lstlisting}
The NVIDIA driver is already on the host --- install a CUDA-matched PyTorch/TF in your env, or just use the container.

%==============================================================================
\section{Storage}
%==============================================================================
\begin{center}
\begin{tabularx}{\linewidth}{@{}l l X@{}}
\toprule
\textbf{Path} & \textbf{Use} & \textbf{Notes} \\
\midrule
\code{\textasciitilde} = \code{/workspace/researchers/<you>} & code + results & \HomeQuota{} quota, backed up nightly \\
\code{/workspace/scratch/<you>} & fast scratch & not backed up; files older than 14 days auto-purged \\
\code{/data/shared} & shared datasets / docs & read-only \\
\code{/data} & bulk (NFS, \textasciitilde{}23\,TB) & ask Kyle for a writable project dir if you need one \\
\bottomrule
\end{tabularx}
\end{center}

%==============================================================================
\section{Moving data}
%==============================================================================
\code{scp} and \code{rsync} use your SSH config, so they route through the tunnel with no extra flags:
\begin{lstlisting}
scp bigfile.tar quarkylab:~/                 # upload
scp quarkylab:~/results.csv .                # download
rsync -avP ./dataset/ quarkylab:~/dataset/   # sync a folder
\end{lstlisting}

%==============================================================================
\section{Long-running work: use tmux}
%==============================================================================
SSH sessions can drop (laptop sleep, network blips). A detached session keeps your job alive:
\begin{lstlisting}
tmux new -s run        # start; launch your training inside
# Ctrl-b then d to detach; reconnect later:
ssh quarkylab
tmux attach -t run
\end{lstlisting}
(\code{screen} works too; \code{nohup cmd \&} for fire-and-forget.)

%==============================================================================
\section{Etiquette}
%==============================================================================
\begin{itemize}
  \item One physical GPU, shared. Running directly (outside SLURM) you are \textbf{not} VRAM-capped --- check \code{nvidia-smi} first and don't OOM others.
  \item For multi-day campaigns there is a \code{research} SLURM partition (interactive allowed, priority over students) if you'd rather queue and preempt-share --- optional, coordinate with Kyle.
\end{itemize}

\vspace{0.6cm}
\begin{center}\small Help: contact Kyle on \Contact.\end{center}

\end{document}
